\documentclass[aspectratio=169,12pt]{beamer}

\usepackage[UTF8]{ctex}
\usepackage{amsmath, amssymb, bm}
\usepackage{booktabs}
\usepackage{array}
\usepackage{mathtools}

\usetheme{Madrid}
\usecolortheme{default}

\title{Estimation: Bias, MSE, Consistency\\從估計量品質到 SPC 基準建立}
\author{黃丞暘}
\institute{統計專論（三）期中報告}
\date{2026/04/21}

\begin{document}

%================================================
\begin{frame}
    \titlepage
\end{frame}

%================================================
\begin{frame}{1. 報告目標與主線}
\textbf{主題：} Estimation --- Bias, MSE, Consistency

\vspace{0.5em}
\textbf{核心問題：}
\begin{itemize}
    \item 當母體參數未知時，我們如何用樣本去估計？
    \item 什麼叫做一個「好的估計量」？
    \item 為什麼估計量的品質會直接影響 SPC 的監控效果？
\end{itemize}

\vspace{0.5em}
\textbf{本報告主線：}
\[
\text{Estimator}
\;\Longrightarrow\;
\text{Bias / Variance / MSE}
\;\Longrightarrow\;
\text{Consistency}
\;\Longrightarrow\;
\text{Phase I estimation in SPC}
\]

\vspace{0.8em}
\textbf{結構：}
\begin{enumerate}
    \item 理論定義
    \item 數學推導
    \item 統計解釋
    \item SPC 應用
\end{enumerate}
\end{frame}

%================================================
\begin{frame}{2. 為什麼「估計」在 SPC 中重要？}
在 SPC 中，製程的 in-control 狀態通常以基準參數描述，例如：
\[
\mu_0,\quad \sigma_0
\]

但實務上這些量通常 \textbf{未知}，必須由 Phase I 歷史資料估計。

\vspace{0.8em}
\textbf{關鍵事實：}
\begin{itemize}
    \item 若 \(\hat{\mu}_0\) 估錯，中心線會偏移
    \item 若 \(\hat{\sigma}_0\) 估錯，控制界限會過寬或過窄
    \item 界限錯誤會導致：
    \begin{itemize}
        \item false alarm 過多
        \item 真正 shift 時偵測太慢
    \end{itemize}
\end{itemize}

\vspace{0.8em}
因此，SPC 不是只要會畫圖，而是先要回答：

\[
\text{How good is the estimator?}
\]
\end{frame}

%================================================
\begin{frame}{3. 估計量與抽樣分配}
設
\[
X_1,\dots,X_n \overset{iid}{\sim} F_\theta
\]
其中 \(\theta\) 是未知參數。

\vspace{0.5em}
\textbf{估計量（estimator）}是樣本的函數：
\[
\hat{\theta}=T(X_1,\dots,X_n)
\]

\vspace{0.8em}
因為 \(X_1,\dots,X_n\) 是隨機的，所以 \(\hat{\theta}\) 也是隨機變數。  
因此我們不能只問「算出多少」，而要問：

\begin{itemize}
    \item \(E(\hat{\theta})\) 是多少？
    \item \(\mathrm{Var}(\hat{\theta})\) 是多少？
    \item \(\hat{\theta}\) 會不會隨 \(n\to\infty\) 趨近真值 \(\theta\)？
\end{itemize}

\vspace{0.8em}
\textbf{這三個問題對應到：}
\[
\text{Bias},\qquad \text{MSE},\qquad \text{Consistency}
\]
\end{frame}

%================================================
\begin{frame}{4. Bias：估計量是否系統性偏離真值？}
\textbf{定義：}
\[
\mathrm{Bias}(\hat{\theta},\theta)=E(\hat{\theta})-\theta
\]

\vspace{0.5em}
若對所有 \(\theta\) 都有
\[
E(\hat{\theta})=\theta,
\]
則稱 \(\hat{\theta}\) 為 \textbf{unbiased estimator}。

\vspace{0.8em}
\textbf{解釋：}
\begin{itemize}
    \item Bias 衡量的是「平均而言估太高還是估太低」
    \item Bias \(>0\)：平均高估
    \item Bias \(<0\)：平均低估
    \item Bias \(=0\)：沒有系統性偏差
\end{itemize}

\vspace{0.8em}
\textbf{例：樣本平均估計母體平均}
若
\[
E(X_i)=\mu,
\quad
\bar X=\frac1n\sum_{i=1}^n X_i,
\]
則
\[
E(\bar X)=\mu
\]
所以 \(\bar X\) 是 \(\mu\) 的 unbiased estimator。
\end{frame}

%================================================
\begin{frame}{5. MSE：只看 unbiased 還不夠}
除了偏差，估計量還有波動程度。

\vspace{0.5em}
\textbf{Mean Squared Error 定義：}
\[
\mathrm{MSE}(\hat{\theta},\theta)=E\big[(\hat{\theta}-\theta)^2\big]
\]

展開可得：
\[
\mathrm{MSE}(\hat{\theta},\theta)
=
\mathrm{Var}(\hat{\theta})
+
\mathrm{Bias}^2(\hat{\theta},\theta)
\]

\vspace{0.8em}
\textbf{推導：}
\[
\hat{\theta}-\theta
=
(\hat{\theta}-E\hat{\theta})+(E\hat{\theta}-\theta)
\]
所以
\[
E[(\hat{\theta}-\theta)^2]
=
E[(\hat{\theta}-E\hat{\theta})^2]
+
(E\hat{\theta}-\theta)^2
\]
因為交叉項期望為 0，因此
\[
\boxed{
\mathrm{MSE}=\mathrm{Var}+\mathrm{Bias}^2
}
\]

\vspace{0.8em}
\textbf{意義：}
一個估計量就算 unbiased，若 variance 很大，仍然可能不好。
\end{frame}

%================================================
\begin{frame}{6. Bias--Variance Tradeoff：簡單例子}
目標：估計 \(\mu\)，假設
\[
X_1,\dots,X_n \sim (\mu,\sigma^2)
\]

考慮兩個估計量：
\[
\hat{\mu}_1=\bar X,
\qquad
\hat{\mu}_2=0.8\bar X
\]

\vspace{0.5em}
\textbf{第一個估計量：}
\[
E(\hat{\mu}_1)=\mu,\qquad
\mathrm{Var}(\hat{\mu}_1)=\frac{\sigma^2}{n}
\]
所以
\[
\mathrm{MSE}(\hat{\mu}_1)=\frac{\sigma^2}{n}
\]

\vspace{0.5em}
\textbf{第二個估計量：}
\[
E(\hat{\mu}_2)=0.8\mu
\]
所以
\[
\mathrm{Bias}(\hat{\mu}_2,\mu)=0.8\mu-\mu=-0.2\mu
\]
且
\[
\mathrm{Var}(\hat{\mu}_2)=0.8^2\mathrm{Var}(\bar X)=0.64\frac{\sigma^2}{n}
\]
因此
\[
\mathrm{MSE}(\hat{\mu}_2)=0.64\frac{\sigma^2}{n}+(0.2\mu)^2
\]

\vspace{0.5em}
\textbf{結論：}
biased estimator 仍可能因 variance 較小而有更低的 MSE。
\end{frame}

%================================================
\begin{frame}{7. 數值例子：biased estimator 可能更好}
令
\[
\sigma^2=4,\qquad n=100,\qquad \mu=0.2
\]

則
\[
\mathrm{MSE}(\hat{\mu}_1)=\frac{4}{100}=0.04
\]

而
\[
\mathrm{MSE}(\hat{\mu}_2)
=
0.64\cdot \frac{4}{100} + (0.2\times 0.2)^2
\]
\[
=
0.0256+0.0016=0.0272
\]

因此
\[
0.0272<0.04
\]

\vspace{0.8em}
\textbf{統計解釋：}
\begin{itemize}
    \item \(\hat{\mu}_1=\bar X\) 沒有 bias，但 variance 較大
    \item \(\hat{\mu}_2=0.8\bar X\) 有一點 bias，但 variance 明顯下降
    \item 當 bias 很小、variance 減少很多時，MSE 反而更小
\end{itemize}

\vspace{0.8em}
\textbf{結論：}
\[
\text{「unbiased」不是唯一標準，MSE 才是更完整的品質衡量。}
\]
\end{frame}

%================================================
\begin{frame}{8. Consistency：樣本變大時，估計會不會靠近真值？}
\textbf{定義：}
若估計量序列 \(\hat{\theta}_n\) 滿足
\[
\hat{\theta}_n \xrightarrow{P} \theta,
\]
則稱 \(\hat{\theta}_n\) 是 \(\theta\) 的 \textbf{consistent estimator}。

\vspace{0.8em}
也就是對任意 \(\varepsilon>0\)，
\[
P(|\hat{\theta}_n-\theta|>\varepsilon)\to 0
\quad (n\to\infty)
\]

\vspace{0.8em}
\textbf{直觀意義：}
\begin{itemize}
    \item 小樣本時，估計量可能誤差很大
    \item 當樣本越來越多，估計量應該越來越穩定地貼近真值
\end{itemize}

\vspace{0.8em}
\textbf{一個常見充分條件：}
若
\[
E(\hat{\theta}_n)\to \theta
\quad\text{且}\quad
\mathrm{Var}(\hat{\theta}_n)\to 0,
\]
則通常可推出 consistency。
\end{frame}

%================================================
\begin{frame}{9. Bernoulli 比例估計：unbiased、UMVUE、consistent}
設
\[
X_1,\dots,X_n \overset{iid}{\sim}\mathrm{Bernoulli}(\pi)
\]
其中
\[
P(X_i=1)=\pi,\qquad P(X_i=0)=1-\pi
\]

考慮樣本比例
\[
\hat p=\bar X=\frac1n\sum_{i=1}^n X_i
\]

\vspace{0.5em}
\textbf{性質：}
\[
E(\hat p)=\pi
\]
所以 \(\hat p\) 是 unbiased estimator。

又因為
\[
\mathrm{Var}(X_i)=\pi(1-\pi),
\]
故
\[
\mathrm{Var}(\hat p)=\frac{\pi(1-\pi)}{n}\to 0
\]
因此
\[
\hat p \xrightarrow{P} \pi
\]
即 \(\hat p\) 是 consistent。

\vspace{0.8em}
\textbf{更進一步：}
\[
T=\sum_{i=1}^n X_i \sim \mathrm{Binomial}(n,\pi)
\]
而 \(\hat p=T/n\) 是 \(T\) 的 unbiased function，因此它是 \(\pi\) 的 UMVUE。
\end{frame}

%================================================
\begin{frame}{10. 從估計走到 SPC：Phase I 的核心}
在 SPC 的 Phase I，我們需要先建立 in-control baseline。

\vspace{0.5em}
以 \(\bar X\)-chart 為例，若已知 \(\mu_0,\sigma_0\)，控制界限可寫成
\[
UCL=\mu_0+z_{1-\alpha/2}\frac{\sigma_0}{\sqrt{m}},
\qquad
CL=\mu_0,
\qquad
LCL=\mu_0-z_{1-\alpha/2}\frac{\sigma_0}{\sqrt{m}}
\]

但實務上通常未知，因此改用估計量：
\[
\hat{\mu}_0,\qquad \hat{\sigma}_0
\]

則實際界限變成
\[
\widehat{UCL}=\hat{\mu}_0+z_{1-\alpha/2}\frac{\hat{\sigma}_0}{\sqrt{m}},
\qquad
\widehat{LCL}=\hat{\mu}_0-z_{1-\alpha/2}\frac{\hat{\sigma}_0}{\sqrt{m}}
\]

\vspace{0.8em}
\textbf{問題：}
\begin{itemize}
    \item 若 \(\hat{\mu}_0\) 有偏，中心線位置會錯
    \item 若 \(\hat{\sigma}_0\) 變異大，界限不穩定
    \item 若估計不一致，樣本再多也無法建立可靠基準
\end{itemize}
\end{frame}

%================================================
\begin{frame}{11. p-chart 的例子：估計量品質如何影響監控}
對屬性資料，若每批樣本大小為 \(n\)，不良率為 \(\pi\)，則
\[
\hat p=\frac{X}{n},\qquad X\sim\mathrm{Binomial}(n,\pi)
\]

在 Phase I 中，通常用
\[
\bar p
\]
作為 \(\pi_0\) 的估計。

對應的 p-chart 控制界限常寫為
\[
UCL=\bar p + z\sqrt{\frac{\bar p(1-\bar p)}{n}},
\qquad
LCL=\bar p - z\sqrt{\frac{\bar p(1-\bar p)}{n}}
\]

\vspace{0.8em}
\textbf{這裡的統計邏輯是：}
\begin{enumerate}
    \item 先用樣本估計 \(\pi_0\)
    \item 再用估計的變異建立控制界限
    \item 最後用新資料判定是否 out-of-control
\end{enumerate}

\vspace{0.8em}
\textbf{因此：}
\[
\text{Estimation quality}
\;\Longrightarrow\;
\text{Control limit quality}
\;\Longrightarrow\;
\text{Monitoring quality}
\]
\end{frame}

%================================================
\begin{frame}{12. 結論}
\textbf{本報告的核心結論有三個：}

\vspace{0.5em}
\textbf{(1) 一個好的估計量不能只看 unbiased}
\[
\mathrm{MSE}=\mathrm{Var}+\mathrm{Bias}^2
\]
因此 estimator 的品質來自偏差與波動的共同作用。

\vspace{0.8em}
\textbf{(2) consistency 是長期可靠性的要求}
\[
\hat{\theta}_n \xrightarrow{P} \theta
\]
表示樣本變大後，估計量應穩定逼近真值。

\vspace{0.8em}
\textbf{(3) 在 SPC 中，估計不是前置小步驟，而是監控的基礎}
\begin{itemize}
    \item Phase I 要估 \(\mu_0,\sigma_0\) 或 \(\pi_0\)
    \item 這些估計決定 control limits
    \item control limits 又決定 false alarm 與 detection power
\end{itemize}

\vspace{0.8em}
\[
\boxed{
\text{若估計基準不可靠，則整個 SPC 監控系統都不可靠。}
}
\]
\end{frame}

%================================================
\begin{frame}{備用頁：可能 Q\&A}
\textbf{Q1. 為什麼 unbiased 不一定最好？}

因為 unbiased 只表示平均不偏，但若 variance 很大，估計仍可能很不穩。  
因此更完整的衡量是
\[
\mathrm{MSE}=\mathrm{Var}+\mathrm{Bias}^2
\]

\vspace{0.8em}
\textbf{Q2. consistency 和 unbiased 有什麼差別？}

\begin{itemize}
    \item unbiased：有限樣本下的平均正確性
    \item consistency：大樣本下是否收斂到真值
\end{itemize}

\vspace{0.8em}
\textbf{Q3. 這和 SPC 最直接的連結是什麼？}

Phase I 必須估計 in-control baseline；估計品質會直接影響 control limits，進而影響監控表現。
\end{frame}

\end{document}